\section{Definitions, normalizations and the modular triple}

\begin{statement}{A.1}{Proved}{The modular dictionary and its polynomial identities}
For the matrices fixed above, the following table defines affine group
schemes over $\Z$. In entries involving solutions of an equation, the
group law is transported through the indicated isomorphism; it is not
asserted to be componentwise multiplication in $(\alpha,\beta)$.
\end{statement}
\begin{longtable}{p{.16\linewidth}p{.56\linewidth}p{.18\linewidth}}
\toprule Symbol&Definition or corrected formula&C3 scan; folio\\\midrule\endhead
$G,SG,PG$&$\GL_2,\SL_2,\PGL_2$&323;462\\
$B_0,B'_0$&Upper triangular invertible matrices;
 $b_0(m,l)=\left(\begin{smallmatrix}m&l\\0&1\end{smallmatrix}\right)$&368;485\\
$T_0,L_0,L_1$&$D(m)=\operatorname{diag}(m,1)$; $L_0(z)$; $L_1(z)$&414;508--509\\
$T_r$&$\{cr+dI:Q_3(c,d)\in R^\times\}$&479;(9)\\
$T_s$&$\{ts+uI:Q_2(t,u)\in R^\times\}$&479;(9)\\
$ST_r,ST_s$&Kernels of the respective determinant (norm) maps&376;489\\
$N'_s$&$\{b\in G:b s b^{-1}=e s,\ e^2=1\}$&325,479\\
$\mathcal X$&$(\zeta,\eta,\mu)$ with determinants $1$,
 $\operatorname{tr}(\zeta r)=1$, $\operatorname{tr}(\eta s)=0$,
 $\zeta=\eta L_1(\mu-1)$&268,370;435,486\\
$\mathcal G_{r,s}$&$[\alpha,r]=[\beta,s]L_1(\mu-1)$,
 $\alpha,\beta\in G$; $\simeq B'_0\times T_r\times T_s$&368;(1)--(2)\\
$\mathcal G^*_{r,s}$&$[\alpha,r]=e[\beta,s]L_1(\mu-1)$,
 $e^2=1$; $\simeq B'_0\times T_r\times N'_s$&371;(9)--(10)\\
$S\mathcal G^*_{r,s}$&Previous row with $(\det\alpha)^2=1$, $\det\beta=1$&373;487\\
$\widetilde S\mathcal M$&$B'_0\times T_r\times N'_s\times\GL_2$ in the coordinates of B.4&393;(43)--(44)\\
$\widetilde{\mathcal M}$&$\Gm\times T_r\times N'_s\times\GL_2$; quotient of the preceding row by $L_0$&393;(43)--(44)\\
$\mathcal H$&$\{(a,b)\in T_r\times N'_s:(\det b/\det a)^2=1\}$&375,416;509\\
$\mathcal M$&$\{(a,b,U):\det U=\det b=e'\det a,\ (e')^2=1\}$&399;(13);479;(3)\\
$S\mathcal M$&$B'_0\times_{\Gm,\det,\chi}\mathcal M$&479;(4)\\
$\chi,e,e'$&$\det U=\det b$; $b s b^{-1}=e s$; $\det b/\det a$&480;(12bis)\\
$\mathcal M^0,\mathcal H^0$&Kernels of $(e,e')$; equal-norm triples or pairs&400,416\\
$\mathcal M(0)$&$(a,b)\mapsto(a,b,D(\det b))$&402--403,417\\
$\mathcal M[0]$&$\mathcal M(0)L_0=\{(a,b,U):U\in B'_0\}$&417--418\\
$\mathcal M(s),\mathcal M(r)$&$U=b$; $U=a$, the latter only over $e'=1$&476,491\\
$\Sigma$&Scalar sign factors in $ST_r\times ST_s\times\SL_2$;
 the scheme is $\mu_2^3$, whereas global discrete signs form $(C_2)^3$&401--402\\
$I^0,J^0$&$\ker(\mathcal M^0\to G)$;
 $\ker(\mathcal M^0\to PG)$&400;(17)--(18)\\\bottomrule
\end{longtable}
\begin{proof}
Solving $Ar=rA$ gives $A=cr+dI$; solving $As=sA$ gives $A=ts+uI$.
Their determinants are the displayed quadratic forms. Solving
$bs=e sb$, with $e^2=1$, gives
\[
 b=\begin{pmatrix}u&t\\-et&eu\end{pmatrix},\qquad
 \det b=e(t^2+u^2).
\]
Thus the sign is unique, including over nonreduced rings. It is a
character because conjugation is multiplicative. The subgroup $N'_s$
is $T_s\rtimes\mu_2$: a splitting is $e\mapsto\operatorname{diag}(1,e)$,
and its action on $ts+uI$ sends $t$ to $et$. All remaining determinant
conditions are equalities between characters and hence define closed
subgroups of the indicated affine groups. Representability of inverses
uses the explicitly inverted norms. The parametrization assertions are
proved in B.1; the subgroup and exact-sequence assertions are proved in
B.4--B.6.

The program \texttt{AB\_identities.py} works over universal polynomial
rings and their indicated localizations. It verifies matrix powers,
both norms, commutator formulas, the parametrizing identity and the
general $(v,w)$ normalization. It does not replace proofs of
representability or exactness by finite testing.
\end{proof}

\begin{statement}{A.2}{Proved}{Exactly when a regular unipotent lift is unique}
Let $R$ be a ring and let $U\in\GL_2(R)$ have trace $2$, determinant $1$
and a non-scalar image in every residue field. The matrices with those
properties and the same image as $U$ in $\PGL_2(R)$ are exactly
\[
 (1+\tau)U,\qquad \tau\in R,\quad 2\tau=0,\quad\tau^2=0.
\]
Here ``same image'' is equality of sections of the fppf quotient;
the scalar relating two global matrix lifts is global. Consequently
uniqueness on $R$-points is equivalent to
$\{\tau\in R:2\tau=\tau^2=0\}=0$. For a scheme $S$ the analogous
statement uses $\Gamma(S,\mathcal O_S)$; uniqueness after every base
change is equivalent to $2$ being invertible on $S$.
\end{statement}
\begin{proof}
If $U'=\lambda U$, determinants give $\lambda^2=1$ and traces give
$2\lambda=2$. With $\tau=\lambda-1$ these equations are precisely
$2\tau=\tau^2=0$. Conversely these equations preserve trace and
determinant, and Cayley--Hamilton gives $(U'-I)^2=0$. They preserve
regularity because $\lambda=1$ in every residue field. The scalar is
global since the kernel sheaf of $\GL_2\to\PGL_2$ is $\Gm$.
If $2$ is invertible the assertion persists under all base changes.
If a characteristic-$2$ point exists, base change to
$k[\delta]/(\delta^2)$ supplies the nonzero $\tau=\delta$.

The ambiguity is represented by the nonflat finite group scheme
\[
 H=\Spec\Z[\tau]/(2\tau,\tau^2),\qquad
 \Delta\tau=\tau\otimes1+1\otimes\tau+\tau\otimes\tau.
\]
This is exactly the manuscript's $\mu'_2$, not a finite flat replacement
for it (C3:358--360, fols.480--481). Over $\Z/4$, the two matrices
$\left(\begin{smallmatrix}1&0\\1&1\end{smallmatrix}\right)$ and
$\left(\begin{smallmatrix}3&0\\3&3\end{smallmatrix}\right)$ are a
counterexample to unrestricted uniqueness. Both rings of cardinality
$2$ and $3$ are fields, so cardinality $4$ is minimal. The computation
checks all stated equalities, not only their reductions modulo $2$.
\end{proof}

\begin{statement}{A.3}{Proved}{Pinning with specified lifts, and its projective ambiguity}
Suppose $R=\Z$, a field, or a local ring, and let $(E,R_0,S_0)$ be a
triple in $\GL_2(R)$ satisfying
\[
 E=-R_0S_0^{-1},\quad \det R_0=\det S_0=1,\quad
 \operatorname{tr}R_0=1,\quad\operatorname{tr}S_0=0.
\]
Assume $E$ is regular unipotent and its fixed line and its inverse
image under $S_0$ are complementary. There is a basis, unique modulo
a common scalar, carrying the triple to
$(L_1(1),r,s_-)$. Therefore the corresponding inner identification
of group schemes is unique. For the pair of projective images alone,
the stabilizer of the standard pair in $\PGL_2(R)$ is
\[
 \{[I+\tau r]:2\tau=\tau^2=0\}.
\]
Thus projective data alone need not determine a unique pinning.
\end{statement}
\begin{proof}
Write $N=E-I$. A regular nilpotent endomorphism of a rank-two free
module has image and kernel the same invertible direct summand $L_1$.
Locally this follows by choosing a unit entry of $N$ and solving
$N^2=0$; the kernel and image constructions then glue. Put
$L_0=S_0^{-1}L_1$. Complementarity implies that
$N:L_0\to L_1$ is an isomorphism. Over the listed rings both lines
are free. Choose $e_0$ generating $L_0$ and put $e_1=N e_0$.
Then $E=L_1(1)$ and $S_0$ has form
$\left(\begin{smallmatrix}0&b\\c&d\end{smallmatrix}\right)$.
Its trace gives $d=0$; from $R_0=-ES_0$ and its trace we get
$b=-1$; the determinant gives $c=1$. Only the common scalar remains.
The same proof for a rank-two vector bundle gives
$V\simeq L\otimes\mathcal O_S^2$, hence an identification of its
endomorphism algebra even when $L$ is not trivial. For an Azumaya
algebra the argument is made fppf locally and descends because the
identification just obtained is unique. This proves the assertion for
specified nilpotent/root-subgroup data on a general base as well.

For projective data, a matrix representative of a conjugator satisfies
$grg^{-1}=\lambda r$. Taking traces gives $\lambda=1$, so
$g=cr+dI$. Write $gs_-=e s_-g$; determinant comparison gives $e^2=1$.
The entries give $c=(e-1)d$ and $(e-1)c=0$. Since $g$ is invertible,
$d$ is a unit, so $\tau=e-1$ satisfies $\tau^2=2\tau=0$ and
$[g]=[I+\tau r]$. Conversely these equations suffice. Over $\Z/4$
the matrix $\left(\begin{smallmatrix}1&2\\2&3\end{smallmatrix}\right)$
is a nonidentity projective stabilizer; this is checked in the code.
For the listed rings every projective conjugator has a matrix
representative because their Picard groups vanish. On a general base
the calculation describes the stabilizer sheaf fppf locally.
The qualification in the margin of C3:323, fol.462, is therefore
essential. The later inclusion of a based additive subgroup in
C3:413, fol.508, repairs the ambiguity.
\end{proof}

\begin{statement}{A.4}{Proved}{Classification of normalized modular triples}
Over each of $R=\Z$, an arbitrary field, $\Z/4$, $\Z/8$ and $\Z/9$,
the triples of A.3 form one $\GL_2(R)$ conjugacy class. Their
$\SL_2(R)$ conjugacy classes are naturally indexed by
$R^\times/(R^\times)^2$. In particular there are respectively
$2,2,4,2$ such classes over $\Z,\Z/4,\Z/8,\Z/9$; for a field the
square-class formula applies in every characteristic, including
imperfect characteristic $2$.
\end{statement}
\begin{proof}
A.3 gives transitivity. The common centralizer of $r,s_-$ is the
scalar group: commuting with their sum $\operatorname{diag}(0,1)$
forces a matrix to be diagonal, and commuting with $s_-$ makes the
diagonal entries equal. Two conjugators yield the same $\SL_2$ orbit
exactly when their determinants differ by the determinant of a
scalar, namely a square. Every unit is the determinant of
$\operatorname{diag}(a,1)$. This proves both completeness and
separation of the invariant.

For finite rings the program enumerates every matrix, every pair with
the prescribed traces and determinants, and the regularity and
transversality conditions in every residue field. Independently it
forms every conjugacy orbit. The results are
\[
\begin{array}{c|rrrrrrr}
R&\F_2&\F_3&\Z/4&\F_5&\F_7&\Z/8&\Z/9\\\hline
\text{triples}&6&24&48&120&336&384&648\\
\GL_2\text{ orbits}&1&1&1&1&1&1&1\\
\SL_2\text{ orbits}&1&2&2&2&2&4&2
\end{array}
\]
Complete representatives and orbit sizes are in
\texttt{AB\_enumeration.json}. This classification concerns the
normalized lifts, not arbitrary matrices merely having powers
$R_0^3=S_0^2=-I$. The latter conditions alone permit additional
triples; transversality and the based unipotent datum cannot be deleted.
\end{proof}

\begin{statement}{A.5}{Proved}{The regular conjugacy condition and the two norm obstructions}
For $z=r$ or $s$, let $\mathcal A_z$ consist of matrices $Y\in\SL_2$
such that $\operatorname{tr}(Yz)=\operatorname{tr}z$, and let
$\mathcal A_z^{\rm reg}$ be the open locus where $Yz$ is non-scalar
in every geometric fibre. The map
\[
 \GL_2\longrightarrow\mathcal A_z^{\rm reg},\qquad
 A\longmapsto[A,z]
\]
is an fppf $T_z$-torsor. The same map from $\SL_2$ is an fppf
$ST_z$-torsor. Over a ring $R$, once a $\GL_2(R)$ conjugator exists,
its $\SL_2(R)$ obstruction is exactly its determinant in
$R^\times/N(\mathcal O_z\otimes R)^\times$.
\end{statement}
\begin{proof}
The trace and determinant conditions say that $Yz$ and $z$ have the
same characteristic polynomial. A matrix that is non-scalar over
every residue field is locally cyclic: a cyclic vector in a residue
field lifts to one after localizing at its cyclic determinant.
Relative to a cyclic basis its matrix is the companion matrix of that
polynomial. Consequently $Yz$ and $z$ are Zariski locally conjugate,
and two conjugators differ on the right by $T_z$. This proves the
first torsor assertion. The norm $T_z\to\Gm$ is fppf surjective:
a unit has a square root after the finite faithfully flat extension
$R[T]/(T^2-a)$, and scalar elements have square norm. Adjusting a
conjugator by such an element gives determinant $1$, proving the
second assertion and the exact point obstruction. A global
$\GL_2(R)$ conjugator can additionally be obstructed by a rank-one
module over the quadratic algebra; local cyclicity does not erase
this Picard obstruction.

For $R=\Z$, the Gaussian and Eisenstein orders are Euclidean for
their norms. Indeed rounding a complex coordinate to the nearest
Gaussian lattice point gives squared distance at most $1/2$; for the
triangular Eisenstein lattice the corresponding bound is $1/3$.
Division with remainder therefore decreases the positive integral
norm. Both ideal class groups vanish. A matrix with either
irreducible characteristic polynomial makes $\Z^2$ a torsion-free
rank-one module over the corresponding order, hence a free one.
Thus global conjugators exist for the numerical targets over $\Z$.
For $r$, the identity
\[
 \det(A)[A,r]_{22}=g^2+gh+h^2
 \quad\left(A=\begin{pmatrix}a&b\\g&h\end{pmatrix}\right)
\]
shows that the $\SL_2(\Z)$ image is exactly $Y_{22}>0$; for $s$ the
identity $[A,s]=AA^t/\det A$ shows it is exactly the positive definite
target, equivalently $Y_{22}>0$. The missed integral points are
therefore specified completely by $Y_{22}<0$ in the respective
numerical target. No finite list is intended for these infinite sets.
\end{proof}

\begin{statement}{A.6}{Refuted}{Specific formula slips and their corrections}
The following formulas, in the stated matrix convention, are false:
the trace condition $B+D-C=1$ in the transcription of C3:266;
$s(b)=b^{-t}$ and $bb^t=eI$ for arbitrary $b\in N'_s$ in C3:480;
and the claim in C3:401 that $ST_s$ is central in the entire signed
$\mathcal M$. Their correct forms are, respectively,
\[
 C+D-B=1,\qquad sbs^{-1}=(\det b)b^{-t},\qquad
 bb^t=e(\det b)I,
\]
and centrality of $ST_r\times ST_s$ in $\mathcal M^0$ only.
\end{statement}
\begin{proof}
Direct multiplication gives $\operatorname{tr}(Yr)=C+D-B$.
The universal $2\times2$ adjugate formula gives
$sbs^{-1}=(\det b)b^{-t}$. The normalizer condition then gives the
displayed equation for $bb^t$. Taking $b=2I$ over $\Q$ disproves
the formula omitting its determinant. Conjugation by the splitting
of $e=-1$ in $N'_s$ inverts $ST_s$. For example $s\in ST_s(\Z)$
is sent to $s^{-1}=-s\ne s$. This is a concrete counterexample to
the asserted centrality. The first correction is supported also by
C3:151's $C+D-B=1$; the incorrect order of $B,C$ was checked
visually on scan266 itself. The centrality restriction is already
corrected in the manuscript at C3:419, (45). The alternatives in C3:328 for the lower-right entry
of $R_0^2S_0$ are tested universally: the correct expression is
$-w(w^2+2v)+v+w^2$. All these checks and the degree-three identities
of C3:327--330 are recorded in the output and readings register.
\end{proof}

\begin{statement}{A.7}{Proved}{Projective modular pairs, including imperfect fields}
Consider projective pairs satisfying the order-$3$, order-$2$,
fppf-local regular unipotence and fibrewise transversality conditions
of C3:323, without specifying a unipotent lift. Over a field $k$ of
characteristic different from $2$ there is one $\GL_2(k)$ orbit,
with $\SL_2(k)$ orbits indexed by $k^\times/k^{\times2}$.
In characteristic $2$, put $K=k[T]/(T^2+T+1)$, a quadratic etale
algebra. The $\GL_2(k)$ orbits are indexed by
\[
 N_{K/k}(K^\times)/k^{\times2}.
\]
Each such orbit has its $\SL_2(k)$ orbits forming a torsor for
$k^\times/k^{\times2}$. Over $\Z,\Z/4,\Z/8,\Z/9$ there is one
$\GL_2$ orbit; the numbers of $\SL_2$ orbits are $2,1,2,2$.
These projective counts differ from A.4's counts of normalized lifts.
\end{statement}
\begin{proof}
Transversality implies that both projective elements are non-scalar
in every residue field. Otherwise the two unipotent fixed lines
would coincide in that fibre. For a representative $R$ of the
order-$3$ element, Cayley--Hamilton gives
\[
 R^3=((\operatorname{tr}R)^2-\det R)R
       -(\operatorname{tr}R)\det R\,I.
\]
Since $R^3$ is scalar and the off-diagonal entries and diagonal
difference generate the unit ideal, $\det R=(\operatorname{tr}R)^2$.
The trace is a unit. Dividing by it gives trace and determinant $1$.
A cyclic basis conjugates it to $r$ over a field or local ring;
over $\Z$ the Eisenstein ideal argument in A.5 gives the same result.
The projective involution has a trace-zero representative $S$, by
the same regularity argument applied to its square.

If $2$ is invertible, a projectively unipotent matrix $V$ has the
unique unipotent normalization $2V/\operatorname{tr}V$.
A.3 applies after scaling $S$ compatibly. In characteristic $2$
the condition is $\operatorname{tr}(rS)=0$, since
$S^{-1}=S/\det S$. The solutions are precisely $S=z s$ with
$z\in K^\times$. Here $s$ induces the nontrivial automorphism
$z\mapsto\bar z$ of $K$ by conjugation, and $\det S=N(z)$.
Scalar changes of $S$ multiply $z$ by $k^\times$; conjugation
preserving $r$ multiplies it by $g/\bar g$ for $g\in K^\times$.
The kernel of the norm consists exactly of such quotients:
if $N(w)=1$ and $w\ne1$, take $g=1+w$, obtaining
$g/\bar g=w$. The split-algebra case obeys the same formula,
and $w=1$ is immediate. Thus the orbit invariant is exactly
the norm modulo squares. These forms are projectively regular
unipotent after adjoining the required square root. If their
two fixed lines coincided geometrically, $r,S$ would preserve
that line; the diagonal entries would then give
$\operatorname{tr}(rS)=\lambda\operatorname{tr}r=\lambda\ne0$.
This checks transversality. A stabilizing $g\in K^\times$ must have
$g/\bar g\in k^\times$. Its norm is $1$, so this scalar has square
$1$ and hence is $1$ in a characteristic-$2$ field. Therefore
$g\in k^\times$. The projective stabilizer on field points is trivial,
and the determinant argument of A.4 proves the $\SL_2$ assertion.

For $R=\Z/2^n$ with $n\ge2$, an fppf unipotent normalization
$\lambda V$ satisfies $\lambda\operatorname{tr}V=2$ and
$\lambda^2\det V=1$. A faithfully flat algebra $B$ preserves
$\operatorname{Ann}_B(2)=2^{n-1}B$. Reduction modulo $2$ first
shows that $\operatorname{tr}V=2a$ in $R$. Then
$\lambda a=1+2^{n-1}b$, so $a$ is a unit and
$\lambda^2a^2=1$. Hence $\det V=a^2$ already in $R$.
The global scalar $a^{-1}$ supplies a unipotent normalization.
Over $\Z$, $(\operatorname{tr}V)^2=4\det V$ with
$\det V\in\{\pm1\}$ likewise gives a global normalization;
over $\Z/9$ the argument with invertible $2$ applies. Scaling
$S$ to make $E=-rS^{-1}$ this normalization makes
$\det S=1$, while its trace stays zero. A.3 now gives the
unique $\GL_2$ orbit. Its stabilizer determinants are
\[
 R^{\times2}\{1+\tau:2\tau=\tau^2=0\}.
\]
Their indices in $R^\times$ give $2,1,2,2$. The finite program
projects all normalized triples and forms their $\SL_2$ orbits,
confirming these different projective counts.

The imperfect-field phenomenon is concrete. Over $\F_2(t)$ take
$S=\left(\begin{smallmatrix}1&t\\t+1&1\end{smallmatrix}\right)$.
Its norm is $t^2+t+1$, whose derivative is $1$ and which is
therefore not a square in $\F_2(t)$. This pair is not conjugate
to the standard projective pair, despite satisfying the weak
projective conditions. The nonsquare certificate is in the code.
\end{proof}

\begin{statement}{A.8}{Refuted}{Unrestricted uniqueness among all group-scheme isomorphisms}
The modular data cannot determine a unique arbitrary group-scheme
isomorphism $\GL_2\to\GL_2$, even over $\Z$ or $\Q$.
The correct uniqueness in A.3 is among algebra-induced inner
identifications. Indeed
\[
 \theta(g)=(\det g)^{-1}g
\]
is a nonidentity involutive automorphism of $\GL_2$ fixing
$\SL_2$ pointwise, hence fixing $r,s$ and both root subgroups.
\end{statement}
\begin{proof}
Determinants are central scalars, so the formula is multiplicative.
In rank two $\det\theta(g)=(\det g)^{-1}$, and hence $\theta^2(g)=g$.
On the centre it sends $aI$ to $a^{-1}I$; an inner or algebra-induced
automorphism fixes that centre pointwise. For example $2I$ goes
to $I/2$ over $\Q$. The program verifies multiplicativity and
involutivity universally. Thus a uniqueness claim encompassing
all group-scheme isomorphisms needs the inner/algebra condition,
independently of the characteristic-$2$ infinitesimal ambiguity.
\end{proof}

\begin{statement}{A.9}{Open}{The general dictionary and its exact remaining comparison}
For the general groups of \S49 we use the following dictionary.
It is a dictionary of discrete or profinite groups, or of sheaves
when those objects are explicitly supplied, and does not silently
assert that their automorphism functors are group schemes.
For the signed rows we impose
$\rho^3=\sigma^2=\omega\in Z(C)$ and $\omega^2=1$,
and retain the inversion involution specified in B.13.
The symbol $-1$ in those rows means this common central element
when $\omega$ has order two.
The signed fibre-product rows additionally use
$L_0^{\rm sgn}\cap Z(C)=1$, as in B.13; the unsigned and signed
cyclic subgroups are distinguished in the table.
\end{statement}
\begin{longtable}{p{.17\linewidth}p{.57\linewidth}p{.16\linewidth}}
\toprule Our symbol&Definition; manuscript symbol&C3 scan\\\midrule\endhead
$C$&The group generated by $\rho,\sigma$; $\mathcal C$ or $\mathcal G$&502;530\\
$L_0,L_1$&Closed cyclic groups generated by $\sigma^{-1}\rho$, $\rho\sigma^{-1}$ in the unsigned convention&502\\
$B$&Automorphisms preserving both conjugacy classes and normalizing $L_0$; $B_{\rho,\sigma}$&503 (7)\\
$\Sigma$&Numerical commutator pairs with $\eta^{-1}\zeta\in L_1$; $\Sigma_{\rho,\sigma}$&503 (8)\\
$\Phi$&$u\mapsto(u(\rho)\rho^{-1},u(\sigma)\sigma^{-1})$; realization is an extra hypothesis&502--503\\
$S_0G,G$&$C\rtimes B$ and its quotient by $(\ell^{-1},\operatorname{Int}\ell)$; $S_0G_{\rho,\sigma},G_{\rho,\sigma}$&505\\
$\gamma$&$B/\operatorname{Int}(L_0)$; $\gamma_{\rho,\sigma}$, previously $\Gamma_{\rho,\sigma}$&506\\
$Z_\rho,Z_\sigma$&Centralizers in $G$; their kernels over $\gamma$ are $C_\rho,C_\sigma$ when $L_0\cap Z(C)=1$&509;511\\
$T^0,ST^0$&$Z_\rho\times_\gamma Z_\sigma$ and $B\times_\gamma T^0$; $\Gamma^0_{\rho,\sigma},S_0\Gamma^0_{\rho,\sigma}$&510--511\\
$M^0,SM^0$&$T^0\times_\gamma G$ and $B\times_\gamma M^0$; $M^0_{\rho,\sigma},S_0M^0_{\rho,\sigma}$&512--514\\
$\widetilde K$&Pullback of a specified global-sign quotient representation; specializes to the definition of $\widetilde N_{1,1}$ after the map is fixed&492;524 (2)--(3)\\
$\varepsilon_2,\varepsilon_3$&Characters of $\widehat\Z^\times$ modulo $4$ and $3$; their additive forms are $\nu_2,\nu_3$&531--534\\
$B^{\rm sgn}$&Pairs $(u,\mu)$ with $u(\sigma^{-1}\rho)=(\sigma^{-1}\rho)^\mu$, $u(\rho)\sim_C\rho^\mu$, $u(\sigma)\sim_C\sigma^\mu$&530 (7)\\
$Z^{\rm sgn}$&Tuples $(u,\mu,e_\rho,e_\sigma)$ with the same loop condition and $u(\rho)\sim_C\rho^{e_\rho}$, $u(\sigma)\sim_C\sigma^{e_\sigma}$; the two signs are independent numerical coordinates&549 (6)\\
$L_0^{\rm sgn},e_1,\ell_1$&$\overline{\langle\sigma\rho\rangle}$, $\rho\sigma$, and $(\rho\sigma)^2$; these are the signed loop conventions used in B.13&530--534\\
$D,\gamma_D,G_D$&$D=B^{\rm sgn}$ or $Z^{\rm sgn}$; $\gamma_D=D/j(L_0^{\rm sgn})$, $G_D=(C\rtimes D)/\{(\ell^{-1},j(\ell))\}$, with $j(\ell)=(\operatorname{Int}\ell,1,1,1)$ and the last signs suppressed for $B^{\rm sgn}$&535--536; 551--553\\
$\Sigma^{\rm eff}$&Numerical signed pairs that are realized by an automorphism; the fourth condition is essential&534--535\\
$N_x,N_x^*,Z_x^*$&For $x=\rho,\sigma$, $N_x=\{(g,e):gxg^{-1}=x^e\}$, $N_x^*=\{(g,e)\in N_x:e=\chi_x(\delta g)\}$, $Z_x^*=C_{G_D}(x)\cap\delta^{-1}\ker\chi_x$, where $\delta:G_D\to\gamma_D$ and $\chi_x$ is the recorded sign&540--541; 555--556\\
$\Gamma_D,S_0\Gamma_D$&$N_\rho^*\times_{\gamma_D}N_\sigma^*$ and $D\times_{\gamma_D}\Gamma_D$; the starred normalizers keep the numerical signs&559 (51), (55)--(56)\\
$\mathcal M_D,S_0\mathcal M_D$&$\Gamma_D\times_{\gamma_D}G_D$ and $D\times_{\gamma_D}\mathcal M_D$; kernels and the cartesian square are proved in B.13&559 (50), (52)--(54); 560 (57), (59)\\
$D_q,H_q,A_q$&For $q:C\twoheadrightarrow C'$, $D_q$ stabilizes $\ker q$; $H_q$ and $A_q$ are the image and kernel of the descended coefficient-group map. Exact formulas are B.15.1--B.15.5&570--572 (88)--(94)\\
$D_0,L(D_0),G_0$&A restricted coefficient subgroup; $L(D_0)=j^{-1}(D_0)$ and $G_0=CD_0\subset G_D$. This keeps all of $C$ as conjugators&571--572 (94)\\
$A,Z^A,L_A,H_A$&The different restriction $A=C^\natural\triangleleft C$ is an open normal conjugating subgroup; $L_A=j^{-1}(Z^A)$ and $H_A=\operatorname{im}(A\rtimes Z^A\to G)$. The antidiagonal uses $L_A\cap A$, as proved in B.17&651--664 (1)--(52 bis)\\
$\mathscr E_{1,1,\Q},\mathscr E'_{1,1,\Q}$&The elliptic moduli stack and its rigidification by $[-1]$, written $M_{1,1},M'_{1,1}$ in the source; their gerbe and coarse $j$-map are distinguished from the affine group $\mathcal M$ in B.16&628 (26)\\
$\mathcal E,\phi,J$&Reconstruction of an extension from $(P,T,J,\theta,\phi)$ satisfying the two compatibilities in B.15.9; compatible gluing additionally requires the diagonal subgroup B.15.11&691--696; fols.~648--650\\\bottomrule
\end{longtable}
\begin{proof}[Proved dictionary, computation, and residual question]
The unsigned entries and every map between them are proved in B.11;
the global-sign pullback is proved in B.9. The finite construction
is executed for $S_3,S_4$, and B.12 supplies a fully executed
failure of $\Sigma^{\rm eff}=\Sigma$ in $\mathrm{PSL}_2(\F_{13})$.
For the signed entries the character definitions follow from
$\rho^3=\sigma^2=\omega$ and $\omega^2=1$.
The generator-image conditions fix $\omega$, because the square
of either conjugate of $\sigma^{\pm1}$ is $\omega$.
Since $\mu$ is odd, replacing $\sigma^{-1}\rho$ by
$\sigma\rho=\omega(\sigma^{-1}\rho)$ then preserves the condition
that it is raised to its $\mu$th power. The two
possible loop conventions must nevertheless not be interchanged
inside unqualified formulas. A formal signed theorem is given in B.13.

The remaining identifications are distinct questions, rather than
one unspecified notation audit. First, numerical signed parameters
must be compared with effective automorphism parameters for
$C=\widehat{\SL_2(\Z)}$.
B.8 defines the numerical triples explicitly and proves that each
gives a continuous endomorphism; its exact remaining assertion
(B.8.5) is surjectivity of every such endomorphism.
The finite unsigned and signed realization computations in B.11
and B.14 are proved cases for that question.

Second, B.15 supplies the quotient functoriality of scans~570--572,
proves stability of the actual congruence kernel, and gives exact
kernel formulas. Its full effective image and compatible lifts
remain separately specified. Third, B.16 treats the central twist
and normalized conjugators of scans~573--582, retaining the exact
named geometric-image and unquotiented-section comparisons of those
scans and scans~612--616. Fourth, the continuation of the dictionary
in B.18 records the normalized sections of scans~583--595, the
normality and index-kernel assertions of scans~635--647, the
bisection of scans~648--650, and the extra open-conjugator datum
of scans~651--690. B.17 refutes the literal smaller quotient and
proves the necessary intersection repair. B.18 specifies the
remaining comparisons first for the named index-six and index-twelve
modular subgroups, with proved formal cases and finite calculations.
The reconstruction and gluing at scans~691--696 are proved in B.15.

D.7 separately records the first modulus-$3$ GT lifting instance
and the remaining arithmetic comparison. Integral representability and comparison
with the schemes of \S48 require an additional construction:
B.8 proves that the naive constant profinite scheme cannot even
realize $\widehat\Z\ni1\mapsto L_0(1)$ over $\Z$.
No untyped fppf quotient or descent theorem is asserted here.
These source-indexed residuals are recorded as partly settled in
A(a), A(b), B(c), and B(g); none is absorbed into (B.8.5).
\end{proof}

